\chapter{Design und Implementierung}
\label{sec:implementierung}

Nachdem sich das vorherige Kapitel mit der Methodik beschäftigt hat, geht dieses Kapitel auf die konkrete Realisierung des Systems ein. Dabei wird erläutert, wie die Architektur aufgebaut ist und wie die einzelnen Softwarekomponenten ineinandergreifen. Darüber hinaus wird auf die verwendeten Datenstrukturen eingegangen.

\section{Softwarearchitektur und Kernkomponenten}
Die Implementierung erfolgte in der Programmiersprache Python und hält sich streng an die Regeln der objektorientierten Programmierung, damit das System wart- und erweiterbar bleibt. Dabei wurde das Kernsystem in drei wesentliche Skripte unterteilt: Erkennungsmodelle (\texttt{models.py}), Zustandsautomat (\texttt{state\_machine.py}) und die Orchestrierung (\texttt{processor.py}).

\subsection{Das Modul \texttt{models.py}: Die \texttt{CatFlapPipeline}}
Die Klasse \texttt{CatFlapPipeline} verwaltet Detektor und Klassifikator und vereinfacht deren Verwendung. Beim Initialisieren des Objekts werden beide Modelle geladen. Dabei werden die Pfade (absolut oder relativ) zu den Modellen ergänzt. Zudem wird angegeben, für welches Modell der CLAHE-Algorithmus angewendet werden soll.

Die Methode \texttt{run\_detector} konvertiert das Bild in Graustufen und verwendet optional den CLAHE-Algorithmus. Anschließend wird die Objekterkennung aufgerufen, durch den Parameter \texttt{persist=True} wird der in Ultralytics integrierte ByteTrack-Tracker aktiviert, welcher uns die Begrenzungsrahmen samt Objekt-IDs liefert.

Die Methode \texttt{run\_classifier} schneidet das Bild wie in Kapitel 4.5 beschrieben zusammen. Im folgenden Codeausschnitt ist zu sehen, wie dieses Zuschneiden implementiert ist und welche Auswirkungen es auf das Katzengesicht hat. Durch das Erweitern nach unten stellen wir sicher, dass das Modell genügend Informationen über mögliche Beute im Maul der Katze erhält.

\begin{PythonCode}[caption={Asymmetrische Randerweiterung für den Bildausschnitt}]
x1, y1, x2, y2 = map(int, box)
# Add padding (heavy on bottom for dangling prey)
pad_w = 0
pad_top = 0
pad_bottom = 30
x1 -= pad_w
y1 -= pad_top
x2 += pad_w
y2 += pad_bottom
\end{PythonCode}
Im nächsten Schritt wird wieder optional der CLAHE-Algorithmus angewendet. Anschließend wird der Bildausschnitt des Katzengesichtes vom Bildklassifikator verarbeitet, dieser liefert uns dann den Konfidenzwert für die Klasse \texttt{CAT\_WITH\_PREY}.

\subsection{Das Modul \texttt{state\_machine.py}: Zeitliche Aggregation}
Der Zustandsautomat ist durch die Klasse \texttt{StateMachine} abgebildet. Der Status einer einzelnen Katze wird in der Klasse \texttt{TrackState} gespeichert, diese speichert die Historie der Konfidenzwerte in einer \texttt{collections.deque} mit einer fixen Maximallänge (\texttt{history\_length = 15}). Durch diese Datenstruktur wird sichergestellt, dass beim Erreichen der Kapazitätsgrenze der älteste Wert aus der Warteschlange entfernt wird  (\glqq First-In-First-Out\grqq). Dabei bleibt das manuelle Speicher-Management komplett erspart.

Der folgende Codeausschnitt zeigt die Kernlogik der \texttt{update}-Methode. Diese berechnet bei jedem neuen Bild den gleitenden Durchschnitt (\texttt{numpy.mean}) und schaltet anschließend den Zustand entsprechend um.

\begin{PythonCode}[caption={Zustandsermittlung mittels gleitendem Durchschnitt}]
def update(self, prey_confidence, frame_idx):
    self.confidence_history.append(prey_confidence)
    self.last_seen = frame_idx

    # Calculate moving average
    avg_confidence = np.mean(self.confidence_history)

    # Determine new state based on aggregated confidence
    if avg_confidence >= self.threshold:
        self.current_state = State.CAT_WITH_PREY
    else:
        self.current_state = State.CAT_NO_PREY
\end{PythonCode}

Die Klasse \texttt{StateMachine} verwaltet die verschiedenen \texttt{TrackState}-Objekte in einem Dictionary (\texttt{track\_id} $\rightarrow$ \texttt{TrackState}). Die Methode \texttt{cleanup\_stale\_tracks} iteriert regelmäßig über dieses Dictionary und entfernt die Objekte, die seit mehr als 30 Einzelbildern nicht mehr vom Tracker zugeordnet wurden konnten, um Speicherlecks zu vermeiden.

Sobald eine oder mehrere Katzen mit Beute im Maul erkannt werden, wirft das \texttt{StateMachine}-Objekt ein Ereignis und meldet, dass der Eintritt von Katzen nicht erlaubt ist. Andernfalls, d. h. wenn alle Katzen sauber sind, meldet das Ereignis, dass der Eintritt erlaubt ist.

\subsection{Die Orchestrierung (\texttt{processor.py})}
Im Skript \texttt{processor.py} geschieht die tatsächliche Verarbeitung der Einzelbilder. Die Klasse \texttt{CatFlapProcessor} verwaltet ein \texttt{StateMachine} Objekt, zählt wie viele Bilder pro Sekunde verarbeitet werden können und bietet die Methode \texttt{process\_frame} an. Diese bildet das Herzstück der Klasse, sie bekommt ein Videoeinzelbild als Parameter übergeben, danach wird über alle erkannten Katzen iteriert, dabei wird für jedes Objekt die Klassifikation durchgeführt. Im letzten Schritt wird dann der Zustandsautomat aktualisiert und die Visualisierungen (Begrenzungsrahmen und Zustand) werden auf das Bild eingezeichnet.

Wie schwierig es ist, qualitativ hochwertige Daten von Katzen mit Beute zu sammeln, wurde im Kapitel 3 bereits erwähnt. Um den Klassifikator gezielt zu verbessern, wurde ein Mechanismus implementiert, der Bilder speichert, bei denen das Modell unsicher war (Konfidenzwert zwischen 0,3 und 0,7). Dadurch können die Szenarien, die dem Modell am schwersten fallen, gezielt verbessert werden. Dieses Feature wird mit dem Kommandozeilenparameter \texttt{--save\_uncertain} aktiviert. Um Speicher zu schonen, wird nach Auslösung des Mechanismus die Speicherung derselben \texttt{track\_id} für 30 Einzelbilder gesperrt.

\section{Die Verbindung zur Hardware (\texttt{catflap.py})}
Die im Skript \texttt{catflap.py} enthaltene Klasse \texttt{Catflap} stellt dabei die Hardwareumgebung zur Verfügung. Durch Polymorphismus können verschiedene Arten erzeugt werden. Dabei liefert die Methode \texttt{capture\_continuous} das nächste Bild der Kamera. Außerdem bietet die Klasse Methoden zum schließen \texttt{lock\_in} / \texttt{lock\_out} und öffnen \texttt{unlock\_in} / \texttt{unlock\_out} der Klappe.

\begin{itemize}
	\item \texttt{lock\_in}: Katze darf Haus nicht betreten.
	\item \texttt{lock\_out}: Katze darf Haus nicht verlassen.
	\item \texttt{unlock\_in} und \texttt{unlock\_out}: Heben die jeweilige Sperre auf.
\end{itemize}

Die Ableitung \texttt{SimulationCatflap} wird zur Simulation verwendet, die Videoeinzelbilder liefert hier ein Video, sperren und entsperren der Klappe wird durch eine Ausgabe auf die Konsole simuliert. Für den realen Betrieb gibt es die Klassen \texttt{ServoCatflap}, die mithilfe eines Servo-Motors sperrt, und \texttt{SureCatflap}, die als Anbindung zu einer vorhandenen SureFlap-Katzenklappe \cite{surepetcare_microchip} verwendet werden kann.

\section{Ausführungsumgebungen}
Das System kann je nach Einsatzzweck in verschiedenen Umgebungen ausgeführt werden, die jeweils eigene Skripte besitzen.

\subsection{Die lokale Testumgebung (\texttt{main.py})}
Das Skript \texttt{main.py} wurde entwickelt um neu trainierte Modelle schnell auf der lokalen Umgebung (z.b. Windows Laptop) validieren zu können. Dabei liest das Skript eine Videodatei ein. Danach durchläuft jedes Videoeinzelbild die \texttt{process\_frame} Methode der Klasse \texttt{CatFlapProcessor}. Anschließend wird die Visualisierung auf das Video eingezeichnet und abgespielt.

\subsection{Die Mikrocontroller Testumgebung (\texttt{pi\_main.py})}
Um das Zusammenspiel aller Komponenten sinnvoll auf dem Endgerät testen zu können wurde das \texttt{pi\_main.py} Skript entwickeltentwickelt. Dadurch ist es möglich, jeden Teil – also Detektor, Klassifikator, Servo und Kamera – isoliert zu testen. Außerdem ist dieser teil wichtig zum zu testen wie viele Bilder pro Sekunde der Raspberry Pi 5 verarbeiten kann. Ansonsten funktioniert das Skript gleich wie \texttt{main.py}.

\subsection{Die Reale Welt (\texttt{pi\_server.py})}
Für den realen Betrieb reichen die vorher erwähnten Skripte nicht aus. Damit man während des Betriebes die Klappe manuell sperren / entsperren und die aktuellen Kamera-Aufnahmen einsehen kann, werden im Skript \texttt{pi\_main.py} mithilfe des \texttt{flask}-Paketes diese Funktionen als HTTP Endpunkte zur Verfügung gestellt.

Auch hier wird ein \texttt{CatFlapProcessor}-Objekt erzeugt und verwaltet, dieses bekommt kontinuierlich die Einzelaufnahmen der Kamera. Sollte sich der Zustand des darin enthaltenen \texttt{TrackState}-Objekts ändern, wird, wie in Kapitel 5.1.2 beschrieben, ein Ereignis geworfen. Sind alle Katzen sauber, wird die Klappe geöffnet, ansonsten wird sie gesperrt. Wurde die Katzenklappe entsperrt, so startet ein Zeitmesser, der die Klappe nach 30 Sekunden wieder schließt.




